% sections/03_method.tex
% §3 Methodology
% §3.1 + §3.2: LR — Nguyễn Công Huy (SE192333)
% §3.3 + §3.4: MS — Lại Minh Dũng (SE194215)
% Nguồn: proposal.md §5.1–§5.5, notes.md, cost_log_LR.md, compute_metrics_MS.py

\section{Methodology}
\label{sec:method}

\subsection{Dataset}
\label{sec:dataset}

We use the Bugzilla bug report dataset released by Acharya and
Ginde~\cite{acharya2025bugquality}, available at
\url{https://github.com/GindeLab/Ease_2025_AI_model}. The
dataset contains structured bug reports mined from the Mozilla
Bugzilla tracker. The raw file
(\texttt{Plus14\_filtered\_bug\_report\_scores\_Summary.xlsx})
comprises 3,966 rows and 32 columns after the authors' initial
quality filtering (requiring S2R, EB, AB, and Additional
Information, and excluding reports containing stack traces or
code snippets).

Each row pairs an \emph{unstructured input} (column
\texttt{NEW\_llama\_output}: a summary generated by Llama3
from the original report) with a \emph{structured ground truth}
(column \texttt{text}: the original Bugzilla report conforming
to the four-section template). No manual annotation was
performed; ground truth is sourced directly from Bugzilla
reporter submissions.

We split the 3,966 rows 80/10/10 (train / validation / test)
using \texttt{random\_state=42}, yielding a test set of 262
samples used for RQ1. The training partition was not used as
we evaluate few-shot prompting only.

\subsection{Experimental Setup}
\label{sec:setup}

\textbf{Model and environment.}
We use \texttt{unsloth/Qwen2.5-7B-Instruct} loaded from
HuggingFace with \texttt{torch\_dtype=float16} and
\texttt{device\_map=``auto''}. Inference runs on a Kaggle
Notebook with a Tesla T4 GPU (free tier, \$0 cost). An
initial attempt with a P100 GPU encountered a CUDA capability
incompatibility (\texttt{sm\_60} not supported); we switched
to T4 before the pilot run.

\textbf{Inference configuration.}
We set \texttt{temperature=0.01} (the minimum permitted value;
Qwen2.5-7B-Instruct requires \texttt{temperature~>~0}),
\texttt{do\_sample=False} (greedy decoding), and
\texttt{max\_new\_tokens=512}.

\textbf{Prompt design.}
We follow the Alpaca-LoRA 3-shot template from
E02~\cite{acharya2025bugquality} (Listing~2), providing
three input--output exemplar pairs drawn from the training
set before the target input. The system prompt instructs
the model to produce four sections: Steps to Reproduce
(S2R), Expected Result (ER), Actual Result (AR), and
Additional Information.

Listing~\ref{lst:prompt} shows the exact template, following
the Alpaca-style format used by E02~\cite{acharya2025bugquality}
(their Listing 2).

\begin{lstlisting}[caption={Alpaca-style prompt template used for 3-shot inference, following E02~\protect\cite{acharya2025bugquality}.}, label={lst:prompt}, basicstyle=\ttfamily\scriptsize, breaklines=true]
You are a senior software engineer specialized in
generating detailed bug reports.

### Instruction:
Please create a bug report that includes the
following sections:
1. Steps to Reproduce (S2R): Detailed steps to
   replicate the issue.
2. Expected Result (ER): What you expected to happen.
3. Actual Result (AR): What actually happened.
4. Additional Information: Include relevant details
   such as software version, build number,
   environment, etc.
If any of these sections are missing from the
provided report, explicitly notify the user which
information is missing.

### Input:
{unstructured_report}

### Response:
{Bug_report}
\end{lstlisting}

The three exemplar pairs (not shown due to space) are
sampled from the training partition; the exact instances
used are available in our replication package.

\textbf{Pilot and full run.}
We first ran a pilot on 80 samples; all 80 outputs contained
the four required sections (100\% valid), confirming pipeline
correctness. We then ran inference on all 262 test samples in
a single Kaggle session with automatic checkpointing every
20 samples. Total runtime was approximately 2.6 hours
($\approx$35 seconds per sample). All 262 outputs were
non-empty and well-formed; 0 samples were excluded. The
replication package includes all outputs in
\texttt{results/full\_llm\_output.csv} and the API log in
\texttt{results/full\_api\_log.txt}.

Our replication package, including datasets, prompts,
scripts, and raw results, is publicly available at:
\url{https://github.com/HoangThong05/RT-SWT-003-nhom3}.

\subsection{Evaluation Metrics}
\label{sec:metrics}

We evaluate generated bug reports using three complementary
metrics, following~\cite{acharya2025bugquality}.

\textbf{CTQRS (Crowdsourced Test Report Quality Score).}
CTQRS~\cite{zhang2022ctqrs} assesses structural completeness
through 13 rule-based sub-checks organized into three
categories: morphological (RM1--4), relational (RR1--5),
and analytical (RA1--4). The maximum possible score is 18
points (RM1--4: $4\times1$; RR1, RR2, RR4, RR5: $4\times1$;
RR3: $1\times2$; RA1--4: $4\times2$). We report
$\text{CTQRS\%} = (\text{raw score} / 18) \times 100$.

We identified and corrected two technical issues in the
GindeLab scorer prior to evaluation: (i) word-list noise in
RA sub-checks causing false positives, and (ii) a hardcoded
\texttt{max\_possible=16} contradicting the actual weighted
sum of 18 derivable from the scorer source code (details in
Section~\ref{sec:threats}). The corrected scorer is
available as \texttt{scripts/ctqrs\_scorer\_v2\_fixed.py}.

\textbf{ROUGE-1 F1.}
ROUGE-1~\cite{lin2004rouge} measures unigram overlap between
generated output and ground-truth reference. We use the
\texttt{rouge-score} library (v0.4.x), reporting F1 score.

\textbf{SBERT Cosine Similarity.}
We compute sentence embedding similarity using
\texttt{sentence-transformers} with model
\texttt{paraphrase-MiniLM-L6-v2}~\cite{acharya2025bugquality},
computing cosine similarity between the generated output and
the ground-truth reference embeddings.

\subsection{Statistical Analysis Plan}
\label{sec:stats}

\textbf{Test selection.}
We use the Wilcoxon signed-rank one-sample
test~\cite{arcuri2011practical} for each of the three
hypotheses in RQ1. This non-parametric test is appropriate
because the metric distributions are not guaranteed to be
normal, and we compare a sample median against a fixed
published threshold. The significance level is $\alpha = 0.05$.

\textbf{Hypotheses.}
For each metric $m \in \{\text{CTQRS}, \text{SBERT},
\text{ROUGE-1}\}$ with threshold $\tau_m$:
$H_0^m$: median$(m) < \tau_m$;
$H_1^m$: median$(m) \geq \tau_m$.
We reject $H_0^m$ when $p < 0.05$ and median$(m) \geq \tau_m$.
Thresholds are taken from Figure~4 of~\cite{acharya2025bugquality}
(GPT-4o 3-shot): CTQRS $\geq 75\%$, SBERT $\geq 0.82$,
ROUGE-1 $\geq 0.47$.

\textbf{Effect size and multiple testing.}
We report Cliff's $\delta$~\cite{arcuri2011practical}
alongside every $p$-value ($|\delta|<0.147$ negligible,
$<0.33$ small, $<0.474$ medium, otherwise large). We apply
Bonferroni correction ($\alpha_{\text{adj}} = 0.05/3
\approx 0.017$) and verify conclusions are unchanged.